% 单调收敛定理证明中的集合 E_n={f_n>=c s}：[0,1]上 f_n(x)=x^{1/n} 递增到 f=1_{(0,1]}；
% 取 s=0.8·1_{(0,1]}，c=0.9，cs=0.72；E_n=[0.72^n,1]：n=1,2,4 时左端点 0.72、0.5184、0.2687。
\begin{tikzpicture}[font=\fontsize{8}{10}\selectfont]
\begin{axis}[width=70mm,height=46mm,xmin=0,xmax=1.12,ymin=0,ymax=1.15,
  axis lines=left,axis line style={FigureStroke,line width=.75pt,-{Stealth[length=4pt]}},
  tick style={FigureStroke,line width=.6pt},tick label style={text=FigureInk},
  xtick={0,0.2687,0.5184,0.72,1},xticklabels={$0$,$0.27$,$0.52$,$0.72$,$1$},ytick={0,0.72,1},yticklabels={$0$,$0.72$,$1$},
  xlabel={$x$},xlabel style={at={(ticklabel* cs:1)},anchor=north west,text=FigureInk},
  clip=false]
  \addplot[FigureStroke,dashed,line width=.6pt,domain=0:1] {0.72};
  \node[anchor=west,text=FigureMuted] at (axis cs:1.01,0.72) {$cs$};
  \addplot[FigureGreen,line width=1.1pt,domain=0:1,samples=100] {x};
  \addplot[FigureYellow,line width=1.1pt,domain=0:1,samples=200] {x^(1/2)};
  \addplot[FigureBlue,line width=1.1pt,domain=0.0001:1,samples=300] {x^(1/4)};
  \node[text=FigureGreen,anchor=north west] at (axis cs:0.6,0.58) {$f_1$};
  \node[text=FigureYellow,anchor=north west] at (axis cs:0.30,0.50) {$f_2$};
  \node[text=FigureBlue,anchor=south east] at (axis cs:0.13,0.86) {$f_4$};
  \foreach \x in {0.2687,0.5184,0.72}{\edef\temp{\noexpand\draw[FigureGrid,line width=.5pt,dashed] (axis cs:\x,0)--(axis cs:\x,0.72);}\temp}
  \addplot[FigureGreen,line width=2pt] coordinates {(0.72,-0.30) (1,-0.30)};
  \addplot[FigureYellow,line width=2pt] coordinates {(0.5184,-0.44) (1,-0.44)};
  \addplot[FigureBlue,line width=2pt] coordinates {(0.2687,-0.58) (1,-0.58)};
  \node[anchor=west,text=FigureInk] at (axis cs:1.02,-0.30) {$E_1$};
  \node[anchor=west,text=FigureInk] at (axis cs:1.02,-0.44) {$E_2$};
  \node[anchor=west,text=FigureInk] at (axis cs:1.02,-0.58) {$E_4$};
\end{axis}
\end{tikzpicture}
